% Nouveau EGA V (cote 25): archive cote 25, site page 20 (PDF page 21).
% typescript. Draft compared: none.
% AI-transcribed by Claude Opus 5.5 (claude-opus-5-5, adaptive thinking, effort high), Message Batches API, batch ega25b (msgbatch_01A7n5YuXubUPNK1zsepi8RG), 2026-11-13; not yet checked by eye.
\scanpage{25/20}
\mspage{\hand{17}}

\uline{Proposition} 4.3.    Soit P l'une des propriétés  suivantes (pour un
préschéma algébrique sur un corps K):

(i)  Z  est lisse sur K.

(ii)  Z satisfait la propriété (R$_k$) géométrique sur K.

(iii)  Z est séparable sur K.

(iv)  Z est géom normal sur K.

\placedmargin{$\uparrow$}(\unsure{vi|v})  Z  est géom. irr.  sur K.

% NOTE: The following overwritten label is circled in pen.
(\unsure{v|vi})  Z est géom. intègre  sur K.

Soit E l'ensemble des \hand{$\xi \in$} $\overset{\hand{\vee}}{P}$  tels que $Y_{\hand{\xi}}$ satisfasse à P. Alors

a)  E est une partie constructible de $\overset{\hand{\vee}}{P}$ .

b)  Supposons  S=Spec(k), k un corps, et supposons dans les cas (i) à (v) que k soit de
car zéro  ou que f:X$\hand{\rightarrow}$P  soit \unsure{un ramifié|non ramifié}. Supposons enfin dans les
cas (v ) et (vi)  que dim $\hand{\overline{\textcolor{black}{f(X)}}}$ $\geqslant$ 2.  Supposons que X satisfasse à P,
alors E contient un ouvert  dense de $\overset{\hand{\vee}}{P}$ .

Démonstration identique à celle de 4.2.   Remarquons que dans
les cas (i) à (\unsure{v}), il suffit de supposer même (dans b) que  f(Z) est
fini, où Z est l'ensemble des points de X où  P est en défaut.

\uline{Corollaire} 4.4.    Considérons aussi la propriété:

($C_m^{}$)   " $\overline{Z}$ ne peut être discon-%
necté par une partie fermée de dimension  $\hand{\leqslant}$ m" (où $\overline{Z}$ est $Z\otimes_K\overline{K}$ , $\overline{K}$
la clôture  algébrique de K).

Soit E l'ensemble des \hand{$\xi \in$} $\overset{\hand{\vee}}{P}$  tels que  $Y_{\hand{\xi}}$ sur k(\hand{$\xi$}) satisfasse à $C_m$.
Alors:

a)  E est constructible.

b)   Supposons S=Spec(k), k un corps, et que  pour toute composante irréduc-%
tible  $X_i$ de X, on a\hand{it} dim $\hand{\overline{\textcolor{black}{f(X_i)}}} \hand{\geqslant} 2$, Alors, si X sur k satisfait $C_{m+1}$ ,
alors E\hand{$_{\xi}$} contient un ouvert dense U de P .

\mspage{\hand{20}}

% NOTE: A pen curve runs upward from the circled lower v/vi label to the preceding item.
